% !TEX program = xelatex
\documentclass[10pt, letterpaper]{article}
\usepackage[UTF8,fontset=none]{ctex}
\setCJKmainfont{Noto Serif CJK TC}
\setCJKsansfont{Noto Sans CJK TC}
\setCJKmonofont{Noto Sans Mono CJK TC}
\usepackage[
    ignoreheadfoot, % set margins without considering header and footer
    top=2 cm, % seperation between body and page edge from the top
    bottom=2 cm, % seperation between body and page edge from the bottom
    left=2 cm, % seperation between body and page edge from the left
    right=2 cm, % seperation between body and page edge from the right
    footskip=1.0 cm, % seperation between body and footer
    % showframe % for debugging 
]{geometry} % for adjusting page geometry
\usepackage[explicit]{titlesec}
\usepackage{tabularx} 
\usepackage{array}
\usepackage[dvipsnames]{xcolor} 
\definecolor{primaryColor}{RGB}{0, 0, 0} 
\definecolor{secondaryColor}{RGB}{0, 0, 170}
\usepackage{enumitem} 
\usepackage{fontawesome5}
\usepackage{amsmath}
\usepackage[
    pdftitle={崔鴻竣履歷},
    pdfauthor={崔鴻竣},
    pdfcreator={LaTeX with RenderCV},
    colorlinks=true,
    urlcolor=secondaryColor
]{hyperref} 
\usepackage[pscoord]{eso-pic}
\usepackage{calc} 
\usepackage{bookmark}
\usepackage{lastpage} 
\usepackage{changepage}
\usepackage{paracol} 
\usepackage{ifthen}
\usepackage{needspace} 
\usepackage{iftex} 
\ifPDFTeX
    \input{glyphtounicode}
    \pdfgentounicode=1
    \usepackage[T1]{fontenc}
    \usepackage[utf8]{inputenc}
    \usepackage{lmodern}
\fi
%\usepackage{newpxtext,newpxmath}
\ifXeTeX
    \setmainfont{TeXGyrePagella}
    \setsansfont{TeXGyreHeros}
    \usepackage{mathpazo}
\else
    \usepackage[osf,sc]{mathpazo}
\fi
\usepackage{euscript} \usepackage{calrsfs}
%%%%%%%%%%%%%%%%%%%%%%%%%%%%%%%%%
\AtBeginEnvironment{adjustwidth}{\partopsep0pt} 
\pagestyle{empty} 
\setcounter{secnumdepth}{0}
\setlength{\parindent}{0pt} 
\setlength{\topskip}{0pt} 
\setlength{\columnsep}{0.15cm}
\makeatletter
\let\ps@customFooterStyle\ps@plain 
\patchcmd{\ps@customFooterStyle}{\thepage}{
    \color{gray}\textit{\small 崔鴻竣－第 \thepage{} 頁／共 \pageref*{LastPage} 頁}
}{}{} 
\makeatother
\pagestyle{customFooterStyle}
%%%%%%%%%%%%%%%%%%%%%%%%%%%%%%%%%
\titleformat{\section}{
    \needspace{4\baselineskip}
    \Large\color{primaryColor}
}{
}{
}{
    \textbf{#1}\hspace{0.15cm}\titlerule[0.8pt]\hspace{-0.1cm}
}[] 

\titlespacing{\section}{
    -1pt
}{
    0.3 cm
}{
    0.2 cm
}
\newcommand{\cvgroup}[1]{%
    \needspace{3\baselineskip}%
    \vspace{0.08cm}%
    {\large\bfseries #1}\par%
    \vspace{0.08cm}%
}
%%%%%%%%%%%%%%%%%%%%%%%%%%%%%%%%%
\newenvironment{one}{
    \begin{adjustwidth}{
        0.2 cm + 0.00001 cm
    }{
        0.2 cm + 0.00001 cm
    }
}{
    \end{adjustwidth}
    \vspace{0.1cm}
} 
%%%%%%%%%%%%%%%%%%%%%%%%%%%%%%%%%
\newenvironment{conf}[2][]{
    \one
    \def\secondColumn{#2}
    \setcolumnwidth{14.5cm, \fill, 5 cm}
    \begin{paracol}{2}
}{
    \switchcolumn \raggedleft \secondColumn
    \end{paracol}
    \endone
    \vspace{0.1cm}
} 
%%%%%%%%%%%%%%%%%%%%%%%%%%%%%%%%%
\newenvironment{pub}[2][]{
    \one
    \def\secondColumn{#2}
    \setcolumnwidth{10.8cm,\fill}
    \begin{paracol}{2}
}{
    \switchcolumn \raggedleft \secondColumn
    \end{paracol}
    \endone
    \vspace{0.1cm}
} 
%%%%%%%%%%%%%%%%%%%%%%%%%%%%%%%%%
\newenvironment{edu}[3][]{
    \one
    \def\thirdColumn{\mbox{#3}}
    \setcolumnwidth{4.2cm,\fill,6.2cm}
    \begin{paracol}{3}
    {\raggedright\mbox{#2}} \switchcolumn
}{
    \switchcolumn \raggedleft \thirdColumn
    \end{paracol}
    \endone
    \vspace{0.1cm}
} 
%%%%%%%%%%%%%%%%%%%%%%%%%%%%%%%%%
\newenvironment{teaching}[3][]{
    \one
    \def\thirdColumn{\mbox{#3}}
    \setcolumnwidth{7.4cm,\fill,3.2cm}
    \begin{paracol}{3}
    {\raggedright\mbox{#2}} \switchcolumn
}{
    \switchcolumn \raggedleft \thirdColumn
    \end{paracol}
    \endone
    \vspace{0.1cm}
} 
%%%%%%%%%%%%%%%%%%%%%%%%%%%%%%%%%
\newenvironment{header}{
    \setlength{\topsep}{0pt}\par\kern\topsep\centering\linespread{1.5}
}{
    \par\kern\topsep
} 
%%%%%%%%%%%%%%%%%%%%%%%%%%%%%%%%%
\newcommand{\CVLastUpdated}{2026/7/31}
\newcommand{\placelastupdatedtext}{% \placetextbox{<horizontal pos>}{<vertical pos>}{<stuff>}
  \AddToShipoutPictureFG*{% Add <stuff> to current page foreground
    \put(
        \LenToUnit{\paperwidth-2 cm-0.2 cm+0.05cm},
        \LenToUnit{\paperheight-1.0 cm}
    ){\vtop{{\null}\makebox[0pt][c]{
        \small\color{gray}\textit{最後更新：\CVLastUpdated}\hspace{\widthof{最後更新：\CVLastUpdated}}
    }}}%
  }%
}%

\let\hrefWithoutArrow\href
\renewcommand{\href}[2]{\hrefWithoutArrow{#1}{\ifthenelse{\equal{#2}{}}{ }{#2 }\raisebox{.15ex}{\footnotesize \faExternalLink*}}}

%%%%%%%%%%%%%%%%%%%%%%%%%%%%%%%%%
\begin{document}
    \newcommand{\AND}{\unskip
        \cleaders\copy\ANDbox\hskip\wd\ANDbox
        \ignorespaces
    }
    \newsavebox\ANDbox
    \sbox\ANDbox{}

    \placelastupdatedtext
%%%%%%%%%%%%%%%%%%%%%%%%%%%%%%%%%
    \begin{header}
        \fontsize{30 pt}{30 pt}
        \textbf{崔鴻竣}

        \vspace{0.3 cm}

        \normalsize
        {\color{secondaryColor}\mbox{{\footnotesize\faMapMarker*}\hspace*{0.13cm} \hrefWithoutArrow{https://www.nthu.edu.tw/}{國立清華大學，臺灣}}}%
        \kern 0.25 cm%
        \AND%
        \kern 0.25 cm%
        \mbox{\hrefWithoutArrow{mailto:hctsui@gapp.nthu.edu.tw}{{\footnotesize\faEnvelope[regular]}\hspace*{0.13cm}hctsui@gapp.nthu.edu.tw}}%
        \kern 0.25 cm%
        \AND%
        \kern 0.25 cm%
        \mbox{\hrefWithoutArrow{https://hctsui.github.io/}{{\footnotesize\faLink}\hspace*{0.13cm}https://hctsui.github.io}}%
    \end{header}
%%%%%%%%%%%%%%%%%%%%%%%%%%%%%%%%%
    \vspace{0.3 cm - 0.3 cm}
%%%%%%%%%%%%%%%%%%%%%%%%%%%%%%%%%
\section{研究領域}
%%%%%%%%%%%%%%%%%%%%%%%%%%%%%%%%%
\begin{one}
函數體上的算術
\end{one}
%%%%%%%%%%%%%%%%%%%%%%%%%%%%%%%%%
\begin{one}
正特徵下的多重 $\zeta$ 值與多重 Eisenstein 級數
\end{one}
%%%%%%%%%%%%%%%%%%%%%%%%%%%%%%%%%
\section{學歷}
%%%%%%%%%%%%%%%%%%%%%%%%%%%%%%%%%
\begin{edu}{\textbf{數學博士班}}{2025 年 2 月至今}
    \textbf{國立清華大學} \\
{\small (逕讀博士)}
\end{edu}
%%%%%%%%%%%%%%%%%%%%%%%%%%%%%%%%%
\begin{edu}{\textbf{數學學士}}{2021 年 9 月至 2024 年 12 月}
    \textbf{國立清華大學} \\
{\small (成績優異畢業)}
\end{edu}
%%%%%%%%%%%%%%%%%%%%%%%%%%%%%%%%%
\section{論文與預印本}
%%%%%%%%%%%%%%%%%%%%%%%%%%%%%%%%%
\cvgroup{預印本}
\begin{pub}{張庭瑋, \underline{崔鴻竣}}
{(1)} \textbf{函數體上 Gross--Koblitz 型公式中 $v$-進 $\Gamma$ 函數的唯一性}\\{\small\color{secondaryColor}[\hrefWithoutArrow{https://drive.google.com/file/d/1-7W2H3vzHCM4E3asVgfFNbzuP1ojudmO/view?usp=drive_link}{PDF} | \hrefWithoutArrow{https://arxiv.org/abs/2504.15697}{arXiv}]}\\\textit{arXiv:2504.15697}, 2025.
\end{pub}
%%%%%%%%%%%%%%%%%%%%%%%%%%%%%%%%%
\begin{pub}{張庭瑋, 陳松筠, 黃斐浚, \underline{崔鴻竣}}
{(2)} \textbf{正特徵任意秩多重 Eisenstein 級數的 $q$-洗牌關係}\\{\small\color{secondaryColor}[\hrefWithoutArrow{https://drive.google.com/file/d/1fzSgTiE4V_rNHtDgoTLtZVMoaWrTHBAx/view?usp=drive_link}{PDF} | \hrefWithoutArrow{https://arxiv.org/abs/2504.18879v3}{arXiv}]}\\\textit{arXiv:2504.18879}, 2025.
\end{pub}
%%%%%%%%%%%%%%%%%%%%%%%%%%%%%%%%%
\begin{pub}{張庭瑋, 陳松筠, 黃斐浚, \underline{崔鴻竣}}
{(3)} \textbf{正特徵多重 Eisenstein 級數的代數結構}\\{\small\color{secondaryColor}[\hrefWithoutArrow{https://drive.google.com/file/d/1odZc8AKcOgvHsn5gvIwKg9MWMdvJ6GcJ/view?usp=sharing}{PDF} | \hrefWithoutArrow{https://arxiv.org/abs/2603.10376}{arXiv}]}\\\textit{arXiv:2603.10376}, 2026.
\end{pub}
%%%%%%%%%%%%%%%%%%%%%%%%%%%%%%%%%
\begin{pub}{\underline{崔鴻竣}}
{(4)} \textbf{正特徵的 $u$-多重 $\zeta$ 值}\\{\small\color{secondaryColor}[\hrefWithoutArrow{https://drive.google.com/file/d/18W5VfIkiHStD1tUsgJRjtaTscF8MsbE6/view?usp=sharing}{PDF} | \hrefWithoutArrow{https://arxiv.org/abs/2604.03618}{arXiv}]}\\\textit{arXiv:2604.03618}, 2026.
\end{pub}
%%%%%%%%%%%%%%%%%%%%%%%%%%%%%%%%%
\section{學術報告}
%%%%%%%%%%%%%%%%%%%%%%%%%%%%%%%%%
\begin{conf}{2026/7/6}
\textbf{正特徵多重 Eisenstein 級數的近期發展}\\{\small\color{secondaryColor}[\hrefWithoutArrow{https://drive.google.com/file/d/1HmrxcOYRdzcAFKnnCzarS6wfa4d7gSbS/view?usp=sharing}{投影片}]}\\東北大學數論研討會, 東北大學, 仙台, 日本
\end{conf}
%%%%%%%%%%%%%%%%%%%%%%%%%%%%%%%%%
\begin{conf}{2026/6/12}
\textbf{正特徵多重 $\zeta$ 值的導子}\\函數體算術工作坊, 國立成功大學, 臺南, 臺灣
\end{conf}
%%%%%%%%%%%%%%%%%%%%%%%%%%%%%%%%%
\begin{conf}{2026/2/8}
\textbf{函數體上的有限多重調和 $u$-級數}\\{\small\color{secondaryColor}[\hrefWithoutArrow{https://drive.google.com/file/d/1lwDJ5mT_b7K84dxeCCNd5DcY-ORmDCHZ/view?usp=sharing}{投影片}]}\\名古屋---關西青年研究者 $\zeta$ 研討會
, 名古屋大學, 名古屋, 日本
\end{conf}
%%%%%%%%%%%%%%%%%%%%%%%%%%%%%%%%%
\begin{conf}{2025/9/30}
\textbf{多重 $\zeta$ 值的 $u$-類比與函數體上 Kaneko--Zagier 猜想的一種可能表述}\\多重 $\zeta$ 值的各種面向, 京都大學, 京都, 日本
\end{conf}
%%%%%%%%%%%%%%%%%%%%%%%%%%%%%%%%%
\begin{conf}{2025/9/3}
\textbf{正特徵多重 Eisenstein 級數的 $q$-洗牌關係}\\{\small\color{secondaryColor}[\hrefWithoutArrow{https://drive.google.com/file/d/1UAwe7PzgwOy-H2U2m3mVj0RLqfC_8Rar/view?usp=sharing}{投影片}]}\\日臺數論會議, 休暇村志賀島, 福岡, 日本
\end{conf}
%%%%%%%%%%%%%%%%%%%%%%%%%%%%%%%%%
\begin{conf}{2025/4/17}
\textbf{Gross--Koblitz--Thakur 公式中 $v$-進 $\Gamma$ 函數的唯一性}\\國立清華大學數論工作坊, 國立清華大學, 新竹, 臺灣
\end{conf}
%%%%%%%%%%%%%%%%%%%%%%%%%%%%%%%%%
\section{獎項與榮譽}
%%%%%%%%%%%%%%%%%%%%%%%%%%%%%%%%%
\begin{conf}{2026}
\textbf{臺日青年科研人才交流計畫}\\公益財團法人日本台灣交流協會
\end{conf}
%%%%%%%%%%%%%%%%%%%%%%%%%%%%%%%%%
\begin{conf}{2025}
\textbf{國科會博士生研究獎學金}\\國家科學及技術委員會
\end{conf}
%%%%%%%%%%%%%%%%%%%%%%%%%%%%%%%%%
\begin{conf}{2024}
\textbf{周鴻經獎學金}\\中央研究院
\end{conf}
%%%%%%%%%%%%%%%%%%%%%%%%%%%%%%%%%
\begin{conf}{2024}
\textbf{中技社數學傑出大學生獎}\\中華民國數學會
\end{conf}
%%%%%%%%%%%%%%%%%%%%%%%%%%%%%%%%%
\section{會議與工作坊}
%%%%%%%%%%%%%%%%%%%%%%%%%%%%%%%%%
\begin{conf}{2026/7/17}
\hrefWithoutArrow{https://sites.google.com/site/kmzsince2011/past/2026?authuser=0}{\textbf{第 72 回關西多重 $\zeta$ 值研討會}}\\東北大學, 仙台, 日本
\end{conf}
%%%%%%%%%%%%%%%%%%%%%%%%%%%%%%%%%
\begin{conf}{2026/7/7 -- 2026/7/10}
\hrefWithoutArrow{https://hiro2.pm.tokushima-u.ac.jp/~hiroki/hiroshima26.html}{\textbf{第 25 屆仙台---廣島數論會議}}\\東北大學, 仙台, 日本
\end{conf}
%%%%%%%%%%%%%%%%%%%%%%%%%%%%%%%%%
\begin{conf}{2026/6/11 -- 2026/6/12}
\hrefWithoutArrow{https://sites.google.com/view/workshoponarithmeticoffunction/home}{\textbf{函數體算術工作坊}} \quad {\color{secondaryColor}\textit{講者}}\\國立成功大學, 臺南, 臺灣
\end{conf}
%%%%%%%%%%%%%%%%%%%%%%%%%%%%%%%%%
\begin{conf}{2026/2/8 -- 2026/2/11}
\hrefWithoutArrow{https://sites.google.com/view/zetayoung19?authuser=0}{\textbf{名古屋---關西青年研究者 $\zeta$ 研討會}} \quad {\color{secondaryColor}\textit{講者}}\\名古屋大學, 名古屋, 日本
\end{conf}
%%%%%%%%%%%%%%%%%%%%%%%%%%%%%%%%%
\begin{conf}{2026/1/21 -- 2026/1/22}
\hrefWithoutArrow{https://2026tms-en.tms.org.tw/}{\textbf{中華民國數學會第 61 屆年會}}\\國立中正大學, 嘉義, 臺灣
\end{conf}
%%%%%%%%%%%%%%%%%%%%%%%%%%%%%%%%%
\begin{conf}{2025/12/15 -- 2025/12/19}
\hrefWithoutArrow{https://www.math.sinica.edu.tw/202512taipeiag}{\textbf{2025 臺北算術幾何會議}}\\中央研究院, 臺北, 臺灣
\end{conf}
%%%%%%%%%%%%%%%%%%%%%%%%%%%%%%%%%
\begin{conf}{2025/9/29 -- 2025/10/3}
\hrefWithoutArrow{https://sites.google.com/view/rimsmzv2025/home}{\textbf{多重 $\zeta$ 值的各種面向}} \quad {\color{secondaryColor}\textit{講者}}\\京都大學, 京都, 日本
\end{conf}
%%%%%%%%%%%%%%%%%%%%%%%%%%%%%%%%%
\begin{conf}{2025/9/1 -- 2025/9/5}
\hrefWithoutArrow{https://www.cck.dendai.ac.jp/math/~namikawa/2025JapanTaiwan/2025JapanTaiwan.html}{\textbf{日臺數論會議}} \quad {\color{secondaryColor}\textit{講者}}\\休暇村志賀島, 福岡, 日本
\end{conf}
%%%%%%%%%%%%%%%%%%%%%%%%%%%%%%%%%
\begin{conf}{2025/8/25 -- 2025/8/29}
\hrefWithoutArrow{https://www.math.ntnu.edu.tw/workshop/eantc2025/}{\textbf{第 11 屆東亞數論會議}}\\國立臺灣大學, 臺北, 臺灣
\end{conf}
%%%%%%%%%%%%%%%%%%%%%%%%%%%%%%%%%
\begin{conf}{2025/6/30 -- 2025/7/4}
\hrefWithoutArrow{https://sites.google.com/ncts.ntu.edu.tw/iwasawa-2025/iwasawa-2025?authuser=0}{\textbf{Iwasawa 2025 國際會議}}\\國立臺灣大學, 臺北, 臺灣
\end{conf}
%%%%%%%%%%%%%%%%%%%%%%%%%%%%%%%%%
\begin{conf}{2025/6/25 -- 2025/6/27}
\hrefWithoutArrow{https://sites.google.com/ncts.ntu.edu.tw/iwasawa-summerschool/}{\textbf{Iwasawa 理論臺北暑期學校}}\\國立臺灣大學, 臺北, 臺灣
\end{conf}
%%%%%%%%%%%%%%%%%%%%%%%%%%%%%%%%%
\begin{conf}{2025/2/17 -- 2025/2/22}
\hrefWithoutArrow{https://www.apu.ac.jp/~rintaro/MK60/indexE.html}{\textbf{第 17 屆日本數學會季節研究所「多重 $\zeta$ 值的發展」第二週}}\\近畿大學, 大阪, 日本
\end{conf}
%%%%%%%%%%%%%%%%%%%%%%%%%%%%%%%%%
\begin{conf}{2025/2/10 -- 2025/2/15}
\hrefWithoutArrow{https://sites.google.com/view/the-17th-msj-siw1}{\textbf{第 17 屆日本數學會季節研究所「多重 $\zeta$ 值的發展」第一週}}\\九州大學西新廣場, 福岡, 日本
\end{conf}
%%%%%%%%%%%%%%%%%%%%%%%%%%%%%%%%%
\begin{conf}{2025/1/14 -- 2025/1/15}
\hrefWithoutArrow{https://2025tms-en.tms.org.tw/}{\textbf{中華民國數學會第 60 屆年會}}\\國立中央大學, 桃園, 臺灣
\end{conf}
%%%%%%%%%%%%%%%%%%%%%%%%%%%%%%%%%
\begin{conf}{2025/1/2 -- 2025/1/3}
\hrefWithoutArrow{https://idv.sinica.edu.tw/mathas/202501SCV/index.html}{\textbf{Nevanlinna 理論、雙曲性與算術動力系統導論}}\\中央研究院, 臺北, 臺灣
\end{conf}
%%%%%%%%%%%%%%%%%%%%%%%%%%%%%%%%%
\begin{conf}{2024/1/22 -- 2024/1/23}
\hrefWithoutArrow{https://2024tms.tms.org.tw/}{\textbf{中華民國數學會第 59 屆年會}}\\國立政治大學, 臺北, 臺灣
\end{conf}
%%%%%%%%%%%%%%%%%%%%%%%%%%%%%%%%%
\begin{conf}{2023/1/16 -- 2023/1/17}
\hrefWithoutArrow{https://www.math.nthu.edu.tw/~2022TMS/home_en.html}{\textbf{中華民國數學會第 58 屆年會}}\\國立清華大學, 新竹, 臺灣
\end{conf}
%%%%%%%%%%%%%%%%%%%%%%%%%%%%%%%%%
\section{教學經歷}
%%%%%%%%%%%%%%%%%%%%%%%%%%%%%%%%%
\cvgroup{國立清華大學}
\begin{teaching}{MATH2410 代數（一）}{2025 秋}
    助教
\end{teaching}
%%%%%%%%%%%%%%%%%%%%%%%%%%%%%%%%%
\begin{teaching}{MATH2420 代數（二）}{2026 春}
    助教
\end{teaching}
%%%%%%%%%%%%%%%%%%%%%%%%%%%%%%%%%
\begin{teaching}{MATH2870 離散數學}{2026 春}
    助教
\end{teaching}
%%%%%%%%%%%%%%%%%%%%%%%%%%%%%%%%%
\begin{teaching}{MATH3450 密碼學}{2025 秋}
    助教
\end{teaching}
%%%%%%%%%%%%%%%%%%%%%%%%%%%%%%%%%
\begin{teaching}{MATH1410 線性代數（一）}{2024 秋}
    助教
\end{teaching}
%%%%%%%%%%%%%%%%%%%%%%%%%%%%%%%%%
\begin{teaching}{MATH1420 線性代數（二）}{2025 春}
    助教
\end{teaching}
%%%%%%%%%%%%%%%%%%%%%%%%%%%%%%%%%
\begin{teaching}{MATH3050 複變數函數論}{2025 春}
    助教
\end{teaching}
%%%%%%%%%%%%%%%%%%%%%%%%%%%%%%%%%
\begin{teaching}{MATH2010 高等微積分（一）}{2023 秋}
    助教
\end{teaching}
%%%%%%%%%%%%%%%%%%%%%%%%%%%%%%%%%
\begin{teaching}{MATH2020 高等微積分（二）}{2024 春}
    助教
\end{teaching}
%%%%%%%%%%%%%%%%%%%%%%%%%%%%%%%%%
\section{個人資訊}
%%%%%%%%%%%%%%%%%%%%%%%%%%%%%%%%%
\begin{one}
    \renewcommand{\arraystretch}{1.2}
    \begin{tabular}{@{} p{3cm} @{\hspace{0.2cm}} p{12cm} @{}}
        \textbf{語言} & 中文（母語） | 英文（CEFR C1） | 日文（CEFR A1） \\
        \textbf{地址} & 臺灣 300044 新竹市東區光復路二段 101 號 \\
        \textbf{電子郵件} & hctsui@gapp.nthu.edu.tw | hctsui.math@gmail.com \\
        \textbf{網站} & \hrefWithoutArrow{https://hctsui.github.io}{https://hctsui.github.io} \\
        \textbf{ORCID} & \hrefWithoutArrow{https://orcid.org/0009-0009-7445-5634}{0009-0009-7445-5634} \\
    \end{tabular}
\end{one}
%%%%%%%%%%%%%%%%%%%%%%%%%%%%%%%%%
\end{document}
